\subsection{Реализация мобильного приложения}

\subsubsection{Выбор программных средств для мобильного приложения}

Выбор технологического стека определялся несколькими требованиями: необходимостью обеспечить работу приложения одновременно на платформах Android и iOS из единой кодовой базы, интеграцией с существующим серверным программным интерфейсом на ASP.NET Core, доступом к аппаратным возможностям устройства (GPS-модуль, камера), а также преемственностью с технологическим стеком диспетчерской веб-панели, реализованной на React и TypeScript.

В качестве основной платформы выбран React Native с языком TypeScript. Главным доводом стала кроссплатформенность и преемственность с веб-панелью диспетчера: подходы к структуре компонентов, управлению состоянием и типизации остаются общими, что снижает порог вхождения и обеспечивает единообразие структурных решений.

Для ускорения разработки использован инструментарий Expo, предоставляющий готовые модули для доступа к камере и геолокации, а также среду Expo Go для тестирования без нативной сборки. Остальные компоненты стека: навигация~--- expo-router с файловой маршрутизацией, доступ к камере~--- expo-image-picker, геолокация~--- expo-location, защищённое хранение токенов~--- expo-secure-store, карта~--- react-native-maps (Google Maps на Android, Apple Maps на iOS), HTTP-взаимодействие~--- Axios с поддержкой JWT-аутентификации.

\subsubsection{Реализация пользовательского интерфейса мобильного приложения}

Пользовательский интерфейс мобильного приложения реализован в соответствии с разработанными макетами. Ниже описаны реализованные экраны.

\textbf{Экран авторизации} (рисунок~\ref{fig:impl-auth}) является первым экраном, отображаемым при запуске приложения для неаутентифицированного пользователя. На экране размещены поля ввода электронной почты и пароля с возможностью переключения видимости символов, ссылка для восстановления пароля, основная кнопка <<Войти>>, альтернативный сценарий входа через внешний аккаунт и переход к регистрации. После успешной аутентификации пользователь перенаправляется на основной экран.

\begin{figure}[H]
    \centering
    \setlength{\fboxrule}{0.4pt}
    \setlength{\fboxsep}{0pt}
    \fbox{\includegraphics[width=0.3\linewidth]{mobile-auth-screen.png}}
    \caption{Экран авторизации мобильного приложения}
    \label{fig:impl-auth}
\end{figure}

\textbf{Экран карты} (рисунок~\ref{fig:impl-map}) является центральным экраном приложения и обеспечивает визуальное отображение обращений на интерактивной карте. Реализация основана на компоненте MapView из библиотеки react-native-maps. Карта отображает маркеры обращений, цвет которых определяется текущим статусом. Реализованы панель фильтров по категориям, кнопка определения текущего местоположения и плавающая кнопка создания обращения.

\begin{figure}[H]
    \centering
    \setlength{\fboxrule}{0.4pt}
    \setlength{\fboxsep}{0pt}
    \fbox{\includegraphics[width=0.35\linewidth]{mobile-map-screen.png}}
    \caption{Экран карты мобильного приложения}
    \label{fig:impl-map}
\end{figure}

\textbf{Экран создания обращения} (рисунок~\ref{fig:impl-create}) предназначен для фиксации проблемы и отправки обращения. Реализован в виде модального окна с формой. Основные компоненты: выпадающий список выбора категории, текстовое поле описания, область прикрепления фотографий с использованием expo-image-picker, интерактивная карта для указания местоположения. При обнаружении аналогичного обращения в радиусе 50~м отображается информационное предупреждение.

\begin{figure}[H]
    \centering
    \setlength{\fboxrule}{0.4pt}
    \setlength{\fboxsep}{0pt}
    \fbox{\includegraphics[width=0.25\linewidth]{mobile-create-report-screen.png}}
    \caption{Экран создания обращения}
    \label{fig:impl-create}
\end{figure}

\textbf{Экран списка обращений} (рисунок~\ref{fig:impl-list}) предназначен для просмотра всех обращений текущего пользователя. Реализован с использованием компонента FlatList. Экран содержит панель вкладок для фильтрации по статусу (<<Все>>, <<Активные>>, <<Завершённые>>), поисковую строку и список карточек. Каждая карточка отображает миниатюру фотографии, категорию и текущий статус.

\begin{figure}[H]
    \centering
    \setlength{\fboxrule}{0.4pt}
    \setlength{\fboxsep}{0pt}
    \fbox{\includegraphics[width=0.25\linewidth]{mobile-reports-list-screen.png}}
    \caption{Экран списка обращений пользователя}
    \label{fig:impl-list}
\end{figure}

\textbf{Экран карточки обращения} (рисунок~\ref{fig:impl-detail}) отображает подробную информацию о выбранном обращении: фотографию проблемы в полную ширину экрана, информационные поля (категория, адрес, дата создания), текстовое описание, счётчик подтверждений от других пользователей и хронологию изменения статусов.

\begin{figure}[H]
    \centering
    \setlength{\fboxrule}{0.4pt}
    \setlength{\fboxsep}{0pt}
    \fbox{\includegraphics[width=0.25\linewidth]{mobile-report-detail-screen.png}}
    \caption{Экран карточки обращения}
    \label{fig:impl-detail}
\end{figure}

\textbf{Экран уведомлений} (рисунок~\ref{fig:impl-notif}) предназначен для информирования пользователя о смене статусов его обращений. Реализован в виде списка уведомлений с иконками типа события, текстовым содержанием и временной меткой. Реализован индикатор непрочитанных уведомлений.

\begin{figure}[H]
    \centering
    \setlength{\fboxrule}{0.4pt}
    \setlength{\fboxsep}{0pt}
    \fbox{\includegraphics[width=0.25\linewidth]{mobile-notifications-screen.png}}
    \caption{Экран уведомлений}
    \label{fig:impl-notif}
\end{figure}

\textbf{Экран профиля} (рисунок~\ref{fig:impl-profile}) предоставляет доступ к личным данным пользователя и настройкам. Включает аватар, имя, блок статистики (общее количество обращений, выполненных и в работе), навигационное меню и кнопку выхода из учётной записи.

\begin{figure}[H]
    \centering
    \setlength{\fboxrule}{0.4pt}
    \setlength{\fboxsep}{0pt}
    \fbox{\includegraphics[width=0.25\linewidth]{mobile-profile-screen.png}}
    \caption{Экран профиля пользователя}
    \label{fig:impl-profile}
\end{figure}

Для обеспечения единообразия интерфейса созданы переиспользуемые компоненты: Button, Input, Card, LoadingIndicator, ErrorDisplay, ReportCard, ImageCarousel. Реализована адаптивность интерфейса с использованием flexbox-разметки, относительных размеров и хука useWindowDimensions.

\subsubsection{Ключевые алгоритмы мобильного приложения}

\textbf{Создание обращения с проверкой дубликатов.} Пользователь заполняет форму: выбирает категорию, вводит описание, прикрепляет фотографию и указывает местоположение. Приложение выполняет первичную валидацию и отправляет запрос POST /api/reports. При обнаружении дубликатов отображается модальное окно с предложением подтвердить существующее обращение или создать новое. При подтверждении отправляется запрос POST /api/reports/\{id\}/confirm.

\textbf{Работа с геолокацией.} При первом обращении к функционалу геолокации приложение запрашивает разрешение через API expo-location. При предоставлении разрешения выполняется запрос getCurrentPositionAsync, координаты передаются для предзаполнения формы и отображения маркера. При отказе предлагается ручное указание на карте.

\textbf{Управление аутентификацией} (рисунок~\ref{fig:alg-auth}). После успешного входа данные сессии сохраняются в защищённом хранилище expo-secure-store. При обращении к защищённым разделам токен автоматически добавляется к HTTP-запросам. При получении ответа 401 выполняется автоматическое обновление токена и повторный запрос. При запуске приложения выполняется попытка восстановления ранее сохранённой сессии.

\begin{figure}[H]
    \centering
    \setlength{\fboxrule}{0.4pt}
    \setlength{\fboxsep}{0pt}
    \fbox{\includegraphics[width=0.35\textwidth]{alg-auth.png}}
    \caption{Алгоритм управления аутентификацией}
    \label{fig:alg-auth}
\end{figure}
